% =====================================================================
% Appendix: Related-work detail.
% CREATED 2026-08-04 (flagship narrative restructure, operation 1). Content
% moved from sections/related_work.tex §"Prior findings on calibration
% sensitivity". NOTHING WAS DELETED IN THE MOVE: the five paragraphs, their
% % SOURCE comments and every qualification appear here verbatim.
%
% The body keeps the one-sentence statement that the two antecedents disagree
% and points here; the reconciliation argument, which is what the objection
% needs answering at length, lives here.
%
% LABEL CHANGE: sec:related:reconcile -> app:reconcile. The two pointing sites
% (sections/related_work.tex, sections/minigrid.tex) were repaired in the same
% commit.
% =====================================================================

\section{Reconciling the calibration-sensitivity antecedents}
\label{app:reconcile}

\textbf{The two most direct antecedents of this paper reach opposite
conclusions about the same knob.} \citet{williamsaletras2024} find
\emph{substantial} variation in downstream task performance across calibration
sets, explicitly contrasting their result with prior work that suggested greater
robustness. \citet{paglieri2024outliers} find the effect \emph{diminishing}:
where the older OPT models degrade badly and vary with the calibration set,
Llama-2 7B, Llama-3 8B, Command-R 35B and Mistral 7B are robust, with Mistral 7B
close to immune, and they argue the field should stop building its quantization
literature around outlier preservation. Both are careful studies of calibration
sensitivity, published within a year of each other, pointing in opposite
directions.

The second of these is also the strongest available objection to this paper,
and it is the one a reader is most likely to arrive already holding: if modern
models are insensitive to the calibration set, why would the calibration seed
reorder anything?

Our answer is that both results hold at once, and that the disagreement between
them is a symptom of measuring the wrong thing for the question the field
actually asks. Three differences do the work, and our own numbers show why.

\paragraph{The intervention is different, and ours is strictly finer.} Both
antecedents vary the calibration \emph{data}. \citet{paglieri2024outliers} do it
across sets differing in quality, content and language: a
RedPajama sample, calibration drawn from uniformly random ASCII punctuation and
whitespace, task-specific sets from ARC-Challenge and PiQA, and multilingual
sets from FLORES+. We hold the corpus fixed and vary only the sample seed, an
intervention they do not make. A finding that swapping the corpus for random
punctuation barely moves accuracy does not entail that redrawing the sample
leaves the \emph{ordering} of two methods intact; the two questions are not
nested in the direction the objection needs.

\paragraph{The unit of analysis is different.} \citet{paglieri2024outliers}
study one method at a time, reporting GPTQ W4A16, AWQ W4A16, SmoothQuant W8A8
and a naive W8A8 baseline in separate result sections, and they report no
head-to-head ranking between methods and no ranking flips. Our unit of analysis
is the \emph{gap between} two methods at the same bit width, a quantity neither
antecedent's design measures.

\paragraph{Individual robustness plus a small method gap implies ranking
instability.} This is the reconciliation proper, and it is what
\S\ref{sec:minigrid} measures. Each method is individually
stable in absolute accuracy, much as they report, and the seed-induced range is
at least as large as the mean GPTQ--AWQ gap in 7 of our 8 confirmatory cells.
% SOURCE: docs/H3_EIGHT_CELL_DECISION_2026-07-26.md (SIGNED), "Mechanical
% application of the rule": range/gap holds in 7 of 8. Wording follows that
% record's own "at least as large as"; do not tighten it to "smaller than".
Those two facts together \emph{are} ranking instability. We are not
contradicting their robustness finding; we are showing it has a consequence they
did not test for. A field that reads ``calibration choice barely moves accuracy''
as licence to compare two methods from one calibration run each has drawn the
wrong inference from a correct result.

\paragraph{The magnitude question.} Robustness claims are
claims about size, so the reconciliation has to be quantitative. The figures
that follow come from a \emph{post-hoc} resolution analysis, prompted by this
very paper and labelled as such wherever it appears (\S\ref{sec:minigrid}); it
is descriptive and modifies no verdict. On MMLU our seed-induced ranges run 5.5
to 17.5 paired standard errors, a spread no robustness claim reaches, and one
the benchmark resolves comfortably at $n = 14{,}042$.
% SOURCE: docs/MINIGRID_SUPPORTING_RESULTS_2026-07-26.md, step 5 table,
% max_range/SE column, MMLU rows: 11.07, 17.45, 5.53, 11.19.
On GSM8K at $n = 1{,}000$ the same ranges run roughly two standard errors, and
we say so instead of averaging the two tasks together: that arm of the
experiment does not resolve the effect it was built to measure
(\S\ref{sec:minigrid:resolution}, \S\ref{sec:limitations}).
% SOURCE: same table, GSM8K rows: 1.92, 2.35, 3.61, 2.45.

\paragraph{What the disagreement was about.} Read this way, the two antecedents
are not really in conflict. \citet{williamsaletras2024} and
\citet{paglieri2024outliers} both measure how much a \emph{single} method's
absolute accuracy moves when the calibration data changes, and they disagree
about the size of that movement on the model generations each studied. Neither
measures what a practitioner comparing two methods actually depends on, namely that the
\emph{ordering} survives, and a quantity nobody measured is one the field has no
basis to assume stable. Our design does not settle their disagreement. It shows
that settling it would not have answered the question either way.

\subsection{Adjacent concurrent measurements}
\label{app:related:concurrent}

% RELOCATED FROM sections/related_work.tex 2026-08-05 (compression pass,
% relocation of extended related work). Moved whole; nothing deleted.
% SOURCE: docs/PRIOR_ART_CONCURRENT_2026-07-24.md §§3-4, both verified against
% the raw arXiv HTML rendering.
\citet{cacioli2026beyondmean} adapts the clinical Reliable Change Index to
per-item comparisons between model \emph{versions} rather than precisions, and
likewise recommends reporting a churn rate beside the mean;
\citet{nikolic2026displacement} extend the flips metric into a leapfrog/drop
decomposition over community-published quantized checkpoints, and find
that KL-divergence proxies lose their ranking signal in precisely the
near-baseline region our certification tables adjudicate.

\subsection{Losslessness definitions and fidelity metrics}
\label{app:related:lossless}

% RELOCATED FROM sections/related_work.tex 2026-08-05.
\citet{helcig2026slq} occupy the phrase this paper is about. They formalise
notions of losslessness for quantized LLMs (task-lossless and the stricter
distribution-lossless), propose the Expected Acceptance Rate as an interpretable
fidelity metric, and ship SLQ, a method that reaches those targets at low bit
widths. Their question and ours are different, and the difference is the whole
of \S\ref{sec:certification}. \emph{They define losslessness and build a method
to achieve it.} \emph{We audit whether existing published claims of losslessness
have the evidence to support them, and compute how many items it would take to
certify one.} Their paper contains no equivalence test at a declared margin, no
power or required-$n$ computation, and no audit of others' claims; ours contains
no quantization method. A practitioner who adopts EAR as a fidelity target still
needs to know how large an evaluation must be before the accuracy half of a
losslessness claim means anything, and that number is what our certification
tables supply.

\subsection{The constructive-audit genre, and preregistration precedent}
\label{app:related:genre}

% RELOCATED FROM sections/related_work.tex 2026-08-05.
The constructive-audit genre is well established: \citet{dodge2019showyourwork}
on reporting the compute and search behind reported numbers, and
\citet{marie2021mtaudit} on the statistical practice of machine-translation
evaluation. We follow both in framing and in tone: the finding is about what the
field reports, not about whether individual authors are right, and the paper's
deliverable is a standard plus the tooling to meet it. The genre has also begun
to adopt preregistration directly, which is the precedent for
\S\ref{sec:prereg}: \citet{gringras2026frontierlag} run a preregistered
bibliometric audit of which models the evaluation literature actually tests, and
\citet{thomas2026prereg} propose preregistering an analysis against a model not
yet released, so that the configuration cannot be tuned against the outcome.

\subsection{The compression families the audited claims span}
\label{app:related:families}

% RELOCATED FROM sections/related_work.tex 2026-08-05.
The audited claims span the two dominant post-training families: weight-only
quantization, represented by GPTQ \citep{gptq2022}, AWQ \citep{awq2023} and
SqueezeLLM \citep{squeezellm2023}; weight-and-activation quantization, by
LLM.int8() \citep{llmint82022} and SmoothQuant \citep{smoothquant2022}; and
one-shot pruning, by SparseGPT \citep{sparsegpt2023} and Wanda
\citep{wanda2023}. We audit the equivalence \emph{language} these papers use
and the evidence offered for it, not the methods themselves, and the audit's
verdicts are about reporting practice, not about whether a method works.
Two of the most aggressive methods in the set are also the most restrained in
what they assert: SpinQuant \citep{spinquant2024} and QuIP\#
\citep{quipsharp2024} report their accuracy deltas without equivalence
language, and are cited here as positive examples alongside the Qwen quantized
model cards and the \texttt{llama.cpp} documentation.
% SOURCE for the honest-non-claimer list:
% docs/RESULTS_2026-07-15_ATLAS_AUDIT.md §2, "Honest non-claimers" bullet.

